\documentclass[a4paper, 14pt]{extarticle}
\usepackage[top=1in, bottom=1in, left=1in, right=1in]{geometry}
\usepackage{amsmath}
\usepackage{amssymb}
\usepackage{graphicx}
\usepackage{fontspec}
\usepackage{hyperref}
\usepackage{tikz}
\usepackage{fontspec}
%\usetikzlibrary{decorations.pathmorphing}
\usetikzlibrary{calc,decorations,patterns,arrows,decorations.pathmorphing,positioning}
\definecolor{pltblue}{HTML}{1F77B4}
% Handwriting font, resolved once into \ppfont and \ppfontname.
%
% "Pretty Neat" (xkcd-style) is the house face and is installed on almost
% nothing. Excalifont is vendored in tools/fonts/ and installed by CI, so it is
% the one face that is always present; Latin Modern is the last resort.
%
% Naming a font that is not installed is a hard error, not a fallback, which is
% why this is a chain and not a line: seven sheets used to name one outright
% and stopped building the day it was uninstalled.
\def\ppfontname{Latin Modern Roman}
\IfFontExistsTF{Pretty Neat}{\def\ppfontname{Pretty Neat}}{%
  \IfFontExistsTF{Humor Sans}{\def\ppfontname{Humor Sans}}{%
    \IfFontExistsTF{Excalifont}{\def\ppfontname{Excalifont}}{}}}
\expandafter\newfontfamily\expandafter\ppfont\expandafter{\ppfontname}
\expandafter\setmainfont\expandafter{\ppfontname}
\tikzset{every picture/.style={/utils/exec={\ppfont}}}


\makeatletter
\pgfset{
  /pgf/decoration/randomness/.initial=2,
  /pgf/decoration/wavelength/.initial=100
}
\pgfdeclaredecoration{sketch}{init}{
  \state{init}[width=0pt,next state=draw,persistent precomputation={
    \pgfmathsetmacro\pgf@lib@dec@sketch@t0
  }]{}
  \state{draw}[width=\pgfdecorationsegmentlength,
  auto corner on length=\pgfdecorationsegmentlength,
  persistent precomputation={
    \pgfmathsetmacro\pgf@lib@dec@sketch@t{mod(\pgf@lib@dec@sketch@t+pow(\pgfkeysvalueof{/pgf/decoration/randomness},rand),\pgfkeysvalueof{/pgf/decoration/wavelength})}
  }]{
    \pgfmathparse{sin(2*\pgf@lib@dec@sketch@t*pi/\pgfkeysvalueof{/pgf/decoration/wavelength} r)}
    \pgfpathlineto{\pgfqpoint{\pgfdecorationsegmentlength}{\pgfmathresult\pgfdecorationsegmentamplitude}}
  }
  \state{final}{}
}
\tikzset{xkcd/.style={decorate,decoration={sketch,segment length=0.5pt,amplitude=0.5pt}}}
\makeatother

\usepackage{etoolbox}
\AtBeginEnvironment{tabular}{\ppfont}

\setlength{\parindent}{0pt}
\setlength{\parskip}{0.5em}
\usepackage{fancyhdr}
\usepackage{geometry}
\usepackage{adjustbox}
\usepackage{titling}
\usepackage{multicol}
\usepackage{amsmath}
\usepackage{amssymb}
\usepackage{graphicx}
\usepackage{hyperref}
\usepackage{tikz}
\usepackage{fontspec}
\usetikzlibrary{calc,decorations,patterns,arrows,decorations.pathmorphing}
\definecolor{pltblue}{HTML}{1F77B4}

\usepackage{amsmath}
\usepackage{tikz}
\usepackage{enumitem}

\title{From Matrix to Map: Network Embedding Worksheet}
\author{Network Science Course}
\date{}

\begin{document}

%\maketitle
\section{Understanding Outer Products}

\begin{enumerate}
    \item Consider two vectors: $\mathbf{a} = [1, 2, 3]$ and $\mathbf{b} = [1,3,1]$. Calculate their outer product $\mathbf{a} \otimes \mathbf{b}$.

    \vspace{3cm}
    \item Without exact calculations, sketch the resulting matrix from the outer product $\mathbf{u} \otimes \mathbf{v}$ of the following pairs of vectors.
    \begin{enumerate}
        \item $\mathbf{u} = [5,4,3,2,1]$ and $\mathbf{v} = [5,4,3,2,1]$
        \item $\mathbf{u} = [1,2,4,0,0]$ and $\mathbf{v} = [3, 1, 2, 0, 0]$
        \item $\mathbf{u} = [1,1,1,-1,-1]$ and $\mathbf{v} = [1,1,1,-1,-1]$
    \end{enumerate}
    Use shading to represent the relative values in the matrix (darker shades for larger values).
    \begin{center}
    \begin{tabular}{ccc}
        (a) & (b) & (c) \\
        \begin{tabular}{|c|c|c|c|c|c|}
        \hline
        & 1 & 2 & 3 & 4 & 5 \\
        \hline
        1 & & & & & \\
        \hline
        2 & & & & & \\
        \hline
        3 & & & & & \\
        \hline
        4 & & & & & \\
        \hline
        5 & & & & & \\
        \hline
        \end{tabular}
        &
        \begin{tabular}{|c|c|c|c|c|c|}
        \hline
        & 1 & 2 & 3 & 4 & 5 \\
        \hline
        1 & & & & & \\
        \hline
        2 & & & & & \\
        \hline
        3 & & & & & \\
        \hline
        4 & & & & & \\
        \hline
        5 & & & & & \\
        \hline
        \end{tabular}
        &
        \begin{tabular}{|c|c|c|c|c|c|}
        \hline
        & 1 & 2 & 3 & 4 & 5 \\
        \hline
        1 & & & & & \\
        \hline
        2 & & & & & \\
        \hline
        3 & & & & & \\
        \hline
        4 & & & & & \\
        \hline
        5 & & & & & \\
        \hline
        \end{tabular}
    \end{tabular}
    \end{center}
\end{enumerate}

\begin{enumerate}[resume]
    \item Look at the matrices you've sketched. If these matrices represented networks, what kind of network structures might each of them represent?
\end{enumerate}

\clearpage
\section{Decomposing Matrices}

\begin{enumerate}[resume]
    \item Consider the following matrix representing a small network:

    \[
    \begin{bmatrix}
    4 & 4 & 3 & 1 & 0 \\
    4 & 4 & 2 & 0 & 1 \\
    3 & 2 & 2 & 0 & 0 \\
    1 & 0 & 0 & 1 & 0 \\
    0 & 1 & 0 & 0 & 1
    \end{bmatrix}
    \]

    Try to ``decompose'' this matrix into the outer product of two vectors with minimal error. Sketch the vectors and use shading to represent their values.
    \vspace{3cm}
    \item Now consider a bipartite network. Rows are 5 students and columns are 4 clubs. Each number is how many times the student went to that club.

    \[
    \begin{bmatrix}
    6 & 3 & 0 & 0 \\
    4 & 2 & 0 & 0 \\
    2 & 1 & 1 & 2 \\
    0 & 0 & 2 & 4 \\
    0 & 0 & 1 & 2
    \end{bmatrix}
    \]

    Can you decompose this into a sum of two outer products? Sketch the vectors for each outer product and use shading.
    \vspace{3cm}
\end{enumerate}

\end{document}